\SourceLeaf{090}
\[
\begin{aligned}
&\log\sin(v^{(0)}-0''583)-\log\sin(v^{(2)}-0''583)\\
&\quad-\log\sin(v^{(3)}-0''382)+\log\sin(v^{(4)}-0''382)\\
&\quad-\log\sin(v^{(6)}-0''414)+\log\sin(v^{(7)}-0''414)\\
&\quad-\log\sin(v^{(16)}-0''389)+\log\sin(v^{(17)}-0''389)\\
&\quad-\log\sin(v^{(19)}-0''368)+\log\sin(v^{(20)}-0''368)=0
\end{aligned}
\]
Superfluum videtur, alteram in forma finita adscribere. His duabus aequationibus respondent sequentes, ubi singuli coëfficientes referuntur ad figuram septimam logarithmorum briggicorum:
\[
\begin{aligned}
&17{,}068(0)-20{,}174(2)-16{,}993(3)+7{,}328(4)\\
&\quad-17{,}976(6)+22{,}672(7)-5{,}028(16)+21{,}780(17)\\
&\quad-19{,}710(19)+11{,}671(20)=-371\\[5pt]
&17{,}976(6)-0{,}880(8)-20{,}617(9)+8{,}564(10)\\
&\quad-19{,}082(13)+4{,}375(14)+6{,}798(18)-11{,}671(20)\\
&\quad+13{,}657(21)-25{,}620(23)-2{,}995(24)+33{,}854(25)=+370
\end{aligned}
\]

Quum nulla ratio indicata sit, cur observationibus pondera inaequalia tribuamus, statuemus \(p^{(0)}=p^{(1)}=p^{(2)}\text{etc.}=1\). Denotatis itaque correlatis aequationum conditionalium eo ordine, quo aequationes ipsis respondentes exhibuimus, per \(A,\allowbreak B,\allowbreak C,\allowbreak D,\allowbreak E,\allowbreak F,\allowbreak G,\allowbreak H,\allowbreak I,\allowbreak K,\allowbreak L,\allowbreak M,\allowbreak N\), prodeunt ad illorum determinationem aequationes sequentes:
\[
\begin{aligned}
-2''197&=5A+C+D+E+H+I+5{,}917N\\
-0{,}436&=6B+E+F+G+I+K+L+2{,}962M\\
-3{,}958&=A+3C-3{,}106M\\
+0{,}722&=A+3D-9{,}665M\\
-0{,}753&=A+B+3E+4{,}696M+17{,}096N\\
+2{,}355&=B+3F-12{,}053N\\
-1{,}201&=B+3G-14{,}707N\\
-0{,}461&=A+3H+16{,}752M\\
+2{,}596&=A+B+3I-8{,}039M-4{,}874N\\
+0{,}043&=B+3K-11{,}963N\\
-0{,}616&=B+3L+30{,}859N\\
-371&=+2{,}962B-3{,}106C-9{,}665D+4{,}696E\\
&\quad+16{,}752H-8{,}039I+2902{,}27M-459{,}33N\\
+370&=+5{,}917A+17{,}096E-12{,}053F-14{,}707G\\
&\quad-4{,}874I-11{,}963K+30{,}859L-459{,}33M+3385{,}96N
\end{aligned}
\]
\SourceLeaf{091}

Hinc eruimus per eliminationem:
\[
\begin{array}{r@{\;}c@{\;}l|r@{\;}c@{\;}l}
A&=&-0{,}598&H&=&+0{,}659\\
B&=&-0{,}255&I&=&+1{,}050\\
C&=&-1{,}234&K&=&+0{,}577\\
D&=&+0{,}086&L&=&-1{,}351\\
E&=&-0{,}477&M&=&-0{,}109792\\
F&=&+1{,}351&N&=&+0{,}119681\\
G&=&+0{,}271&&&
\end{array}
\]
Denique errores maxime plausibiles prodeunt per formulas
\[
\begin{aligned}
(0)&=C+17{,}068M\\
(1)&=A+C\\
(2)&=C-20{,}174M\\
(3)&=D-16{,}993M
\end{aligned}
\]
etc., unde obtinemus valores numericos sequentes; in gratiam comparationis apponimus (mutatis signis) correctiones a clar. \textsc{de Krayenhof} observationibus applicatas:
\[
\begin{array}{r@{\;}c@{\;}r|r||r@{\;}c@{\;}r|r}
&&&\text{\textsc{de Kr.}}&&&&\text{\textsc{de Kr.}}\\[4pt]
(0)&=&-3''108&-2''090&(14)&=&+0''795&+2''400\\
(1)&=&-1{,}832&+0{,}116&(15)&=&+0{,}061&+1{,}273\\
(2)&=&+0{,}981&-1{,}982&(16)&=&+1{,}211&+5{,}945\\
(3)&=&+1{,}952&+1{,}722&(17)&=&-1{,}732&-7{,}674\\
(4)&=&-0.719&+2{,}848&(18)&=&+1{,}265&+1{,}876\\
(5)&=&-0{,}512&-3{,}848&(19)&=&+2{,}959&+6{,}251\\
(6)&=&+3{,}648&-0{,}137&(20)&=&-1{,}628&-5{,}530\\
(7)&=&-3{,}221&+1{,}000&(21)&=&+2{,}211&+3{,}486\\
(8)&=&-1{,}180&-1{,}614&(22)&=&+0{,}322&-3{,}454\\
(9)&=&-1{,}116&0&(23)&=&-2{,}489&0\\
(10)&=&+2{,}376&+5.928&(24)&=&-1{,}709&+0{,}400\\
(11)&=&+1{,}096&-3{,}570&(25)&=&+2{,}701&+2.054\\
(12)&=&+0{,}016&+2{,}414&(26)&=&-1{,}606&-3{,}077\\
(13)&=&-2{,}013&-6{,}014&&&&
\end{array}
\]
